\section{Experimental Setup}
\label{sec:setup}

We design experiments to answer three research questions mapping to the three-step Goodhart saturation mechanism:

\textbf{RQ1:} Is there a statistically significant structural break in PwC benchmark residual CoV at some paper\_count threshold? \textit{(H-E1)}

\textbf{RQ2:} Does the pre-breakpoint benchmark segment exhibit substantially higher performance variance than the global baseline? \textit{(H-M1)}

\textbf{RQ3:} Does the post-breakpoint segment exhibit substantially lower variance than the pre-breakpoint segment? \textit{(H-M2, H-M3)}

\subsection{Dataset}

PwC benchmark archive, August 2026 snapshot, $N=115$ benchmarks after $\text{min\_papers}=38$ filtering, $\text{paper\_count} \in [38, 352]$.

\subsection{Baselines}

\begin{table}[h]
\centering
\caption{Baselines and their roles.}
\label{tab:baselines}
\begin{tabular}{lp{4cm}p{4cm}}
\toprule
Method & Description & Why Included \\
\midrule
OLS trend & Linear regression CoV $\sim$ paper\_count ($\rho=0.137$) & Tests null hypothesis: saturation is a smooth trend with no structural break \\
$S_{\text{index}}$~\cite{Akhtar2026} & Composite LLM saturation metric & Closest existing saturation index --- we contrast our threshold vs.\ composite approach \\
Liao et al.\ aggregate CoV~\cite{Liao2022} & Time-based saturation metric & Largest saturation study --- we contrast paper\_count threshold vs.\ time-based aggregate \\
\bottomrule
\end{tabular}
\end{table}

\subsection{Evaluation Metrics}

\textbf{RQ1:} Permutation $p$ (primary, gate $< 0.05$), paper\_count$^*$ location (gate $\in [10, 120]$), piecewise $F$-test $p$ (corroborative).

\textbf{RQ2:} Pre-segment variance ratio $> 1.0$; $F$-test one-tailed $p < 0.10$.

\textbf{RQ3:} Brown-Forsythe variance ratio $< 1.0$, $p < 0.05$ (H-M2); $\geq$ 2/4 directional metrics at $p < 0.10$ (H-M3).
